\documentclass[11pt,a4paper]{article}
\usepackage[margin=1.6cm]{geometry}
\usepackage[utf8]{inputenc}
\usepackage[T1]{fontenc}
\usepackage{lmodern}
\usepackage[italian]{babel}
\usepackage{xcolor}
\usepackage{tikz}
\usetikzlibrary{arrows.meta,positioning,calc,shapes.geometric,patterns}
\usepackage{booktabs,array,tabularx}
\usepackage{tcolorbox}
\usepackage{enumitem}
\usepackage{siunitx}
\sisetup{output-decimal-marker={,},detect-all=true}
\usepackage{caption}
\pagestyle{plain}

\definecolor{SapienzaRed}{RGB}{130,35,45}
\definecolor{SoftGray}{RGB}{245,245,245}
\definecolor{DarkBlue}{RGB}{25,70,130}
\definecolor{DarkGreen}{RGB}{20,110,75}
\definecolor{AccentOrange}{RGB}{220,120,30}

\newtcolorbox{datibox}{colback=SoftGray,colframe=SapienzaRed,arc=2mm,boxrule=0.8pt,left=5pt,right=5pt,top=5pt,bottom=5pt}
\newcommand{\titolo}[1]{\begin{center}\sffamily\bfseries\Large\color{SapienzaRed}#1\end{center}\vspace{2mm}}
\newcommand{\sottotitolo}[1]{\vspace{1mm}\noindent{\sffamily\bfseries\color{DarkBlue}#1}\par\vspace{1mm}}

\begin{document}

\titolo{Esercizi da foto - Processi e Impianti}
\begin{center}
\small Testo ricostruito dalle due immagini caricate, con schemi ridisegnati in \LaTeX/TikZ.
\end{center}
\vspace{3mm}

\section*{1. Esercizio di distillazione}

\begin{datibox}
Una miscela di composti organici con una portata massica totale di \SI{5500}{kg/h} viene immessa in una colonna di distillazione a piatti come \textbf{liquido saturo} per essere separata nei suoi due componenti. La percentuale in peso del composto più volatile è del \textbf{15\%}. Si vuole ottenere un distillato con una frazione molare del composto più volatile pari a \textbf{0,90}. I pesi molecolari dei due composti sono rispettivamente \textbf{85} e \textbf{160}. Il rapporto di riflusso ottimale è fissato a \textbf{1,5 volte} il rapporto di riflusso minimo.
\end{datibox}

\sottotitolo{Calcolare}
\begin{enumerate}[label=\arabic*.]
\item Le portate orarie di \textbf{distillato} e di \textbf{corpo di fondo} in uscita dalla colonna.
\item Il \textbf{numero di piatti teorici} richiesto per la separazione con il metodo grafico di McCabe-Thiele, avendo a disposizione la relativa curva di equilibrio.
\item Considerando lo schema di processo riportato in Figura 1, in cui un recupero termico viene effettuato usando il corpo di fondo \(B\) come liquido saturo per preriscaldare l'alimentazione \(F\) in colonna, calcolare la \textbf{temperatura del prodotto di fondo} dopo lo scambio termico. L'alimentazione prima dello scambiatore si trova a \SI{20}{\celsius}.
\end{enumerate}

\sottotitolo{Dati termici}
\begin{center}
\renewcommand{\arraystretch}{1.25}
\begin{tabularx}{0.92\textwidth}{lcc}
\toprule
\textbf{Corrente} & \textbf{Temperatura di ebollizione (\si{\celsius})} & \textbf{Calore specifico (\si{kJ/(kmol.\celsius)})} \\
\midrule
Alimentazione & 156 & 60 \\
Distillato & 117 & 55 \\
Corpo di fondo & 169 & 85 \\
\bottomrule
\end{tabularx}
\end{center}

\begin{figure}[h!]
\centering
\begin{tikzpicture}[scale=0.9,>=Latex,thick]
  % column
  \draw[rounded corners=10pt, fill=gray!8] (5,0) rectangle (6.1,5.4);
  \foreach \y in {0.55,0.95,...,4.8}{\draw[gray!60] (5.08,\y) -- (6.02,\y);}
  \draw (5.55,5.4) -- (5.55,5.95) -- (7.2,5.95);
  \draw[->] (7.2,5.95) -- (8.3,5.95) node[right]{\large \(D\)};
  \draw (7.15,5.55) rectangle (8.0,5.75);
  \draw (7.55,6.05) circle (0.13);
  \draw (7.43,5.93) -- (7.68,6.18);

  % feed and heat exchanger
  \draw[->] (1.0,2.6) -- (3.1,2.6) node[midway,above]{\large \(F\)};
  \draw (3.1,2.6) circle (0.22);
  \draw (2.95,2.45) -- (3.25,2.75);
  \draw[->] (3.32,2.6) -- (5.0,2.6) node[midway,above]{\large \(F\)};

  % bottom line and reboiler
  \draw (5.55,0) -- (5.55,-0.55) -- (3.1,-0.55) -- (3.1,2.38);
  \draw[->] (5.55,-0.55) -- (5.55,-1.0) node[below]{\large \(B\)};
  \draw (6.1,0.6) -- (7.2,0.6) -- (7.2,1.5) -- (6.1,1.5);
  \draw (7.2,1.05) circle (0.18);
  \draw (7.05,0.9) -- (7.35,1.2);

  % recuperated B to exchanger and outlet
  \draw[->] (3.1,3.6) -- (1.1,3.6) node[left]{\large \(B\)};
  \draw (3.1,3.6) -- (3.1,2.82);
  \draw (5.55,-0.55) -- (3.1,-0.55);

  \node[align=center,font=\small] at (5.55,-1.75){Colonna a piatti};
\end{tikzpicture}
\caption{Schema di processo semplificato della colonna di distillazione con recupero termico.}
\end{figure}

\newpage
\section*{2. Esercizio di estrazione liquido-liquido}

\begin{datibox}
Si vuole purificare, attraverso un'\textbf{estrazione con solvente in corrente incrociata}, una miscela costituita da un solvente \(B\) e da un soluto \(A\), utilizzando un solvente \(C\) parzialmente miscibile con \(B\). La miscela binaria di partenza \(L_0\) ha una portata massica pari a \SI{180}{kg/h} e una composizione in \(A\), come frazione massica, pari a \textbf{0,40}. Si vuole arrivare a una miscela binaria purificata con frazione massica di \(A\) pari a \textbf{0,10}. La portata ottimizzata di solvente estraente \(V_0\) è pari a \SI{120}{kg/h}. Utilizzare il diagramma di fase ternario riportato in Figura 2.
\end{datibox}

\sottotitolo{Calcolare}
\begin{enumerate}[label=\arabic*.]
\item Il \textbf{numero di stadi estrattivi} necessari alla separazione.
\item Le \textbf{portate orarie di raffinato ed estratto} in uscita dall'ultimo stadio.
\item La \textbf{composizione dell'estratto e del raffinato} in uscita dal primo stadio estrattivo.
\item La \textbf{composizione del composto A} nelle miscele binarie di estratto e raffinato in uscita dal terzo stadio estrattivo.
\end{enumerate}

\begin{figure}[h!]
\centering
\begin{tikzpicture}[scale=10,>=Latex]
  \coordinate (B) at (0,0);
  \coordinate (A) at (1,0);
  \coordinate (C) at (0,1);

  % triangle
  \draw[very thick] (B) -- (A) -- (C) -- cycle;
  \node[left] at (B) {\Large B};
  \node[right] at (A) {\Large A};
  \node[above] at (C) {\Large C};

  % axes ticks bottom and left
  \foreach \x in {0,0.1,...,1.0}{
    \draw[thin] (\x,0) -- +(0,-0.012);
    \draw[gray!25] (\x,0) -- (0,{1-\x});
  }
  \foreach \y in {0,0.1,...,1.0}{
    \draw[thin] (0,\y) -- +(-0.012,0);
    \draw[gray!25] (0,\y) -- ({1-\y},\y);
  }
  \node[font=\scriptsize,below] at (0,-0.025) {0};
  \node[font=\scriptsize,below] at (0.1,-0.025) {0,1};
  \node[font=\scriptsize,below] at (0.3,-0.025) {0,3};
  \node[font=\scriptsize,below] at (0.4,-0.025) {0,4};
  \node[font=\scriptsize,below] at (0.5,-0.025) {0,5};
  \node[font=\scriptsize,below] at (0.6,-0.025) {0,6};
  \node[font=\scriptsize,below] at (0.7,-0.025) {0,7};
  \node[font=\scriptsize,below] at (0.8,-0.025) {0,8};
  \node[font=\scriptsize,below] at (0.9,-0.025) {0,9};
  \node[font=\scriptsize,below] at (1,-0.025) {1};
  \node[font=\scriptsize,left] at (-0.025,0) {0};
  \foreach \y/\lab in {0.1/{0,1},0.2/{0,2},0.3/{0,3},0.4/{0,4},0.5/{0,5},0.6/{0,6},0.7/{0,7},0.8/{0,8},0.9/{0,9},1/{1}}{
    \node[font=\scriptsize,left] at (-0.025,\y) {\lab};
  }

  % approximate binodal curve from photo
  \draw[very thick,DarkGreen]
    plot[smooth] coordinates {
      (0.00,0.02) (0.18,0.025) (0.38,0.035) (0.56,0.055) (0.66,0.10) (0.69,0.18) (0.66,0.28) (0.58,0.40) (0.46,0.54) (0.32,0.67) (0.16,0.80) (0.03,0.88)
    };

  % tie lines approximated
  \foreach \p/\q in {0.03/0.88,0.08/0.82,0.13/0.76,0.18/0.70,0.23/0.64,0.28/0.58,0.33/0.52,0.38/0.46,0.43/0.39,0.48/0.32,0.53/0.25,0.58/0.18,0.62/0.11}{
    \draw[gray!65] (\p,0.025) -- (\q,{1-\q});
  }

  % clean labels
  \node[font=\small,DarkGreen,align=center] at (0.47,0.13) {curva\ di\ lacuna};
  \node[font=\small,align=center] at (0.55,0.72) {diagramma ternario\\ schematico};
\end{tikzpicture}
\caption{Diagramma ternario schematico ricostruito dalla foto. Le linee interne rappresentano le tie-lines usate per la lettura grafica.}
\end{figure}

\vfill
\begin{center}
\small\color{gray!70} Documento generato automaticamente da testo e figure ricostruiti dalle immagini.
\end{center}

\end{document}
